%!TEX root = ../main.tex
%!TEX program = xelatex
% Chapter 2: Words 251 - 500
\chapter{Everyday Communication}
\label{chap:chapter2}

{\small\sffamily\color{mutedtext} Essential vocabulary for daily activities, family, time, places, and common conversational expressions.\par}

\vspace{0.4em}
\markboth{Chapter 2: Everyday Communication}{Words 251--500}

\begin{multicols}{2}
\vocentry{251}{Бывать}{by'vat'}{To be at, to visit}{Жаль, что я не могу бывать у бабушки чаще.}{I wish I could be at my granny's more often.}
\vocentry{252}{Полный}{'polnyj}{Full, complete}{Это ещё не полный потенциал оборудования.}{It's not yet the full potential of the equipment.}
\vocentry{253}{Прийти}{prij'ti}{To come to}{Они извинились за то, что не смогут прийти на похороны.}{They apologized for not being able to come to the funeral.}
\vocentry{254}{Палец}{'palets}{A finger; a toe}{Повар порезал себе палец пока чистил картофель.}{The cook cut his finger while he was peeling potatoes.}
\vocentry{255}{Россия}{ra'sija}{Russia}{Россия богата природными ресурсами.}{Russia is rich in natural resources.}
\vocentry{256}{Любой}{lu'boj}{Any}{Любой другой человек испугался бы, но не мой отец.}{Any other person would get scared but not my father.}
\vocentry{257}{История}{is'torija}{History; a story}{История хранит много секретов.}{History keeps lots of secrets.}
\vocentry{258}{Наконец}{naka'njets}{At last, in the end}{Наконец, стороны пришли к соглашению.}{At last the parties reached an agreement.}
\vocentry{259}{Мысль}{'mysl'}{A thought}{Это неожиданная, но очень полезная мысль.}{It's an unexpected but a very useful thought.}
\vocentry{260}{Узнать}{uz'nat'}{To recognize; to find out}{Я не смог узнать её из-за тёмных очков.}{I couldn't recognize her because of the dark glasses.}
\vocentry{261}{Назад}{na'zad}{Back, backwards}{Путешественникам не хотелось поворачивать назад.}{The travellers didn't want to turn back.}
\vocentry{262}{Общий}{'obstchij}{Common, public}{У них есть общий недостаток – они слишком много разговаривают.}{They have a common drawback – they talk too much.}
\vocentry{263}{Заметить}{za'metit'}{To notice}{Дети не могли не заметить, что мама чем-то расстроена.}{The children couldn't help but notice that their mother was upset with something.}
\vocentry{264}{Словно}{'slovnə}{Like, as if}{Я так хорошо себя чувствую, словно стала на пять лет моложе.}{I feel so good, like I'm five years younger.}
\vocentry{265}{Прошлый}{'proshlyj}{Last, past}{В прошлый раз мы отослали слишком мало приглашений.}{Last time we sent too few invitations.}
\vocentry{266}{Уйти}{uj'ti}{To leave, to go away}{Очень грубо с твоей стороны уйти не попрощавшись.}{It's very rude of you to leave without saying goodbye.}
\vocentry{267}{Известный}{iz'vesnyj}{Known, familiar}{Хорошо известный писатель дал интервью читателям.}{A well-known writer gave an interview to the readers.}
\vocentry{268}{Давно}{dav'no}{Long ago}{Интересно, как давно они купили новую машину.}{I wonder how long ago they bought a new car.}
\vocentry{269}{Слышать}{'slyshat'}{To hear}{Родители говорили громко, и дети могли слышать каждое слово.}{The parents were talking loudly and the children could hear every word.}
\vocentry{270}{Слушать}{'slushat'}{To listen to}{Я не могу больше слушать эту чепуху!}{I can't listen to this nonsense anymore!}
\vocentry{271}{Бояться}{ba'jatsa}{To fear, to be afraid of}{Не нужно бояться плохих снов.}{One shouldn't fear bad dreams.}
\vocentry{272}{Сын}{syn}{A son}{Мой сын окончил университет в этом году.}{My son has graduated from University this year.}
\vocentry{273}{Нельзя}{nelzj'a}{One cannot, it's prohibited}{Здесь нельзя пользоваться мобильным.}{One cannot use a mobile here.}
\vocentry{274}{Прямо}{pr'jamə}{Straight; openly}{Группа шла прямо, а у почты повернула направо.}{The group was going straight ahead and turned right at the post office.}
\vocentry{275}{Долго}{'dolgə}{Long (adv)}{Я хочу, чтобы мои родители жили долго.}{I want my parents to live long.}
\vocentry{276}{Быстрый}{'bystryj}{Quick, fast}{Мне нужен быстрый ответ.}{I need a quick answer.}
\vocentry{277}{Лес}{'les}{Forest, woods}{Осенью мы часто ходим в лес за грибами.}{In autumn we often go to the forest to pick up mushrooms.}
\vocentry{278}{Похожий}{pa'hozhyj}{Similar, alike}{В детстве у меня был похожий кукольный домик.}{In my childhood I had a similar doll house.}
\vocentry{279}{Пора}{pa'ra}{It is time; a season}{Пора принять окончательное решение.}{It is time to make a final decision.}
\vocentry{280}{Пять}{p'jat’}{Five}{Сегодня моей племяннице исполнилось пять.}{My niece turned five today.}
\vocentry{281}{Глядеть}{g'ljadet’}{To look, to watch}{Собака могла часами глядеть из окна, ожидая возвращения хозяина.}{The dog could look out of the window for hours waiting for the master’s return.}
\vocentry{282}{Оно}{a'no}{It (third-person neuter singular pronoun)}{Не ешь это яблоко – оно гнилое.}{Don’t eat this apple – it is rotten.}
\vocentry{283}{Сесть}{'sest}{To sit down}{В этой комнате негде сесть.}{There’s nowhere to sit down in this room.}
\vocentry{284}{Имя}{'imja}{A name}{Это имя очень необычное для нашей страны.}{This name is very unusual for our country.}
\vocentry{285}{Ж}{zh}{A particle that makes an objection by pointing to an obvious fact - after all, but (A shorter form of ‘же’)}{Почему ты мне не веришь? Я ж говорю правду!}{Why don’t you believe me? After all I’m telling the truth!}
\vocentry{286}{Разговор}{razga'vor}{A talk, a conversation}{У меня был очень неприятный разговор с начальством.}{I had a very unpleasant talk with the management.}
\vocentry{287}{Тело}{'telə}{A body}{Её тело просто божественное!}{Her body is just divine!}
\vocentry{288}{Молодой}{mala'doj}{Young}{Кто этот молодой человек на фото?}{Who’s this young man in the photo?}
\vocentry{289}{Стена}{ste'na}{A wall}{Великая Китайская стена – это уникальное сооружение.}{The Great Chinese Wall is a unique construction.}
\vocentry{290}{Красный}{'krasnyj}{Red}{Красный – это цвет страсти.}{Red is the color of passion.}
\vocentry{291}{Читать}{tchi'tat’}{To read}{Мой младший брат может читать часами.}{My younger brother can read for hours.}
\vocentry{292}{Право}{'pravə}{Right; law}{Каждый гражданин имеет право голосовать.}{Every citizen has the right to vote.}
\vocentry{293}{Старик}{sta'rik}{An old man}{«Старик и море» – одно из лучших произведений Хемингуэя.}{“The Old Man and the Sea” is one of the best works by Hemingway.}
\vocentry{294}{Ранний}{'rannij}{Early (Adj)}{У нас был спокойный ранний завтрак.}{We had a peaceful early breakfast.}
\vocentry{295}{Хотеться}{ha'tetsa}{To want (impersonal)}{Тебе должно хотеться большего от жизни.}{You must want more from life.}
\vocentry{296}{Мама}{'mama}{Mama, mummy}{Его первым словом было «мама».}{His first word was “mama”.}
\vocentry{297}{Оставаться}{asta'vatsa}{To remain, to stay}{Безопасность будет оставаться твоей обязанностью.}{Security will remain your responsibility.}
\vocentry{298}{Высокий}{'vysokij}{Tall, high}{Он не достаточно высокий, чтобы играть в баскетбол.}{He’s not tall enough to play basketball.}
\vocentry{299}{Путь}{put’}{Way, path}{Тебе нужно найти свой путь развития.}{You should find your own way of development.}
\vocentry{300}{Поэтому}{pa'ɛtəmu}{Therefore, that is why}{Я простудился, поэтому не поехал в горы.}{I caught a cold, therefore I didn’t go to the mountains.}
\vocentry{301}{Совершенно}{saver'shennə}{Absolutely, completely}{Этот прибор совершенно бесполезен.}{This device is absolutely useless.}
\vocentry{302}{Кроме}{'krome}{Except; besides}{Все, кроме меня, сдали тест.}{Everyone, except me, passed the test.}
\vocentry{303}{Тысяча}{'tysjacha}{A thousand}{У неё была тысяча причин бросить работу.}{She had a thousand reasons to quit her job.}
\vocentry{304}{Месяц}{'mesjats}{A month; moon}{Декабрь – это зимний месяц.}{December is a winter month.}
\vocentry{305}{Брать}{brat’}{To take}{Я не советую тебе брать кредит.}{I don’t recommend you to take a loan.}
\vocentry{306}{Написать}{napi'sat’}{To write}{Автору удалось написать книгу за несколько месяцев.}{The author managed to write the book in a few months.}
\vocentry{307}{Целый}{'tselyj}{Whole; entire}{Ученый потратил целый год на это изобретение.}{The scientist spent a whole year on this invention.}
\vocentry{308}{Огромный}{ag'romnyj}{Huge, vast}{У этой технологии огромный потенциал.}{This technology has a huge potential.}
\vocentry{309}{Начинать}{natchi'nat’}{To start, to begin}{Мы не будем начинать конференцию без него.}{We won’t start the conference without him.}
\vocentry{310}{Спина}{spi'na}{A back}{После вчерашней работы в саду у меня сильно болит спина.}{My back hurts badly after yesterday’s work in the garden.}
\vocentry{311}{Настоящий}{nasta'jastchij}{Real, true; present}{Этот пожарный – настоящий герой.}{This firefighter is a real hero.}
\vocentry{312}{Пусть}{pust’}{Let + subject + the verb}{Пусть все твои дни будут счастливыми!}{Let all your days be happy!}
\vocentry{313}{Язык}{ja'zyk}{Language; tongue}{Нужно много терпения, чтобы выучить иностранный язык.}{A lot of patience is needed to learn a foreign language.}
\vocentry{314}{Точно}{'tochnə}{Exactly, definitely}{Я точно знаю, что делать.}{I know exactly what to do.}
\vocentry{315}{Среди}{sre'di}{Among}{Он самый опытный работник среди своих коллег.}{He’s the most experienced worker among his colleagues.}
\vocentry{316}{Чувствовать}{'tchustvəvət’}{To feel}{Ты не должна чувствовать себя виноватой в том, что произошло.}{You shouldn’t feel guilty for what’s happened.}
\vocentry{317}{Сердце}{'sertse}{Heart}{Это предательство разбило его сердце.}{This betrayal broke his heart.}
\vocentry{318}{Вести}{ves'ti}{To lead}{Я уверен, что эта дорога не может вести в город.}{I’m sure this road can’t lead to the city.}
\vocentry{319}{Иногда}{inag'da}{Sometimes}{Иногда моя сестра ведёт себя достаточно эгоистично.}{My sister behaves quite selfishly sometimes.}
\vocentry{320}{Мальчик}{'mal’chik}{A boy}{Мой племянник очень активный мальчик.}{My nephew is a very active boy.}
\vocentry{321}{Успеть}{us'pet’}{To have time, to make it, to be in time for}{Мы можем не успеть увидеть все достопримечательности.}{We may not have time to see all the sights.}
\vocentry{322}{Небо}{'nebə}{Sky}{Смотреть на звёздное небо – это очень романтично.}{Watching the starry sky is very romantic.}
\vocentry{323}{Живой}{zhi'voj}{Live, alive, living}{Это был живой концерт классической музыки.}{That was a live concert of classical music.}
\vocentry{324}{Смерть}{'smert’}{Death}{Смерть президента стала шоком для всей страны.}{The president’s death came as a shock for the whole country.}
\vocentry{325}{Продолжать}{pradal'zhat’}{To continue, to go on}{Я буду продолжать помогать этой семье.}{I’ll continue to help this family.}
\vocentry{326}{Девушка}{'devushka}{A girl}{Это девушка моей мечты.}{It’s the girl of my dreams.}
\vocentry{327}{Образ}{'obraz}{Image; style}{Образ матери был центральным в его книгах.}{The image of a mother was the central one in his books.}
\vocentry{328}{Ко}{ko}{To; towards (variant of ‘к’ used before consonants)}{Отношение профессора ко мне очень изменилось.}{The professor’s attitude to me changed a lot.}
\vocentry{329}{Забыть}{za'byt’}{To forget}{Как ты мог забыть о годовщине нашей свадьбы?}{How could you forget about our wedding anniversary?}
\vocentry{330}{Вокруг}{vak'rug}{Around, round}{Всё вокруг было покрыто снегом.}{Everything around was covered in snow.}
\vocentry{331}{Письмо}{pis’'mo}{A letter; writing}{Почтальон потерял моё письмо.}{The postman lost my letter.}
\vocentry{332}{Власть}{'vlast’}{Power, authority}{У королевы только номинальная власть.}{The Queen only has nominal power.}
\vocentry{333}{Чёрный}{'tchjornyj}{Black}{Чёрный считается цветом деловой одежды.}{Black is considered to be the color of business clothes.}
\vocentry{334}{Пройти}{praj'ti}{To pass, to pass by}{Должен пройти по крайней мере год, прежде чем ситуация изменится.}{At least a year should pass before the situation changes.}
\vocentry{335}{Появиться}{paja'vitsa}{To appear; to show up}{Это жирное пятно не могло появиться здесь само.}{This greasy stain couldn’t appear here by itself.}
\vocentry{336}{Воздух}{'vozduh}{Air}{Заводы загрязняют воздух.}{Factories contaminate the air.}
\vocentry{337}{Разный}{'raznyj}{Different; varied}{У них разный цвет лица.}{They have different complexion.}
\vocentry{338}{Выходить}{vyho'dit’}{To go out}{Я не рекомендую тебе выходить из дома – там идёт сильный дождь.}{I don’t recommend you to go out of the house – it’s raining heavily there.}
\vocentry{339}{Просить}{pra'sit’}{To ask; to request}{Капитан был слишком гордым, чтобы просить о помощи.}{The captain was too proud to ask for help.}
\vocentry{340}{Брат}{brat}{A brother}{Его брат будет моим соперником в завтрашнем матче.}{His brother will be my opponent in tomorrow’s match.}
\vocentry{341}{Собственный}{'sobstvenyj}{One’s own}{Она открыла свой собственный бизнес.}{She started her own business.}
\vocentry{342}{Отношение}{atna'shenije}{Attitude}{Его отношение к учёбе заслуживает уважения.}{His attitude to studies deserves respect.}
\vocentry{343}{Затем}{za'tem}{Then}{На первом свидании мы сходили в кафе, а затем прогулялись по парку.}{On the first date we went to a cafe and then walked through the park.}
\vocentry{344}{Пытаться}{py'tatsa}{To try}{Не нужно пытаться обмануть меня – я уже всё знаю.}{You shouldn’t try to deceive me – I know everything already.}
\vocentry{345}{Показать}{paka'zat’}{To show}{Мне не терпится показать тебе, что я купил.}{I can’t wait to show you what I bought.}
\vocentry{346}{Вспомнить}{vs'pomnit’}{To recollect, to remember}{Бабушка не могла вспомнить, где видела этого человека.}{Grandmother couldn’t recollect where she’d seen that man.}
\vocentry{347}{Система}{sis'tema}{A system}{Ни одна политическая система не идеальна.}{No political system is perfect.}
\vocentry{348}{Четыре}{tche'tyre}{Four}{Мы звонили им четыре раза – ответа нет.}{We called them four times – no answer.}
\vocentry{349}{Квартира}{kvar'tira}{Apartment, flat}{Наша квартира находится на пятом этаже.}{Our apartment is on the fifth floor.}
\vocentry{350}{Держать}{der'zhat’}{To hold; to keep}{Я не могу больше держать эту тяжёлую сумку.}{I can’t hold this heavy bag anymore.}
\vocentry{351}{Также}{'takzhe}{Also, too, as well}{Я также хочу обратить ваше внимание на проблемы медицины.}{I also want to pay your attention to the problems of medicine.}
\vocentry{352}{Любовь}{lju'bov’}{Love}{Я не уверен, что это любовь.}{I’m not sure it’s love.}
\vocentry{353}{Солдат}{sal'dat}{A soldier}{Любой солдат скучает по дому.}{Any soldier misses his home.}
\vocentry{354}{Откуда}{at'kuda}{Where from}{Откуда эти апельсины?}{Where are these oranges from?}
\vocentry{355}{Чтоб}{tchtob}{So that, in order to, so as to (a short variant of ‘чтобы’)}{Я оставила тебе записку, чтоб ты знал, где я.}{I left you a note so that you know where I am.}
\vocentry{356}{Называть}{nazy'vat’}{To call; to name}{Я предпочитаю называть людей по имени.}{I prefer to call people by the name.}
\vocentry{357}{Третий}{'tretij}{Third}{Третий вопрос оказался самым простым.}{The third question turned out to be the easiest one.}
\vocentry{358}{Хозяин}{ha'zjain}{An owner; a master}{Кто хозяин этой роскошной квартиры?}{Who’s the owner of this luxury flat?}
\vocentry{359}{Вроде}{'vrode}{Like, such as, kind of}{Люди вроде него раздражают меня.}{People like him irritate me.}
\vocentry{360}{Уходить}{uha'dit’}{To leave}{Нам совсем не хотелось уходить с вечеринки.}{We didn’t want to leave the party at all.}
\vocentry{361}{Подойти}{padaj'ti}{To come to, to approach}{Не мог бы ты подойти ближе ко мне – я плохо слышу.}{Could you come closer to me – I don't hear well.}
\vocentry{362}{Поднять}{padn'jat’}{To raise}{Учитель попросил учеников поднять руки, если они согласны с высказыванием.}{The teacher asked the students to raise their hands if they agreed with the expression.}
\vocentry{363}{Особенно}{a'sobennə}{Especially}{На приёме мне особенно понравилась музыка.}{At the reception I especially liked the music.}
\vocentry{364}{Спрашивать}{sp'rashivat’}{To ask about}{Очень невежливо спрашивать о личных вопросах.}{It’s very impolite to ask about personal matters.}
\vocentry{365}{Начальник}{na'chal’nik}{A boss, a chief}{Наш начальник на удивление понимающий человек.}{Our boss is a surprisingly understanding person.}
\vocentry{366}{Оба}{'oba}{Both (for masculine nouns)}{Оба примера отлично подходят для описания ситуации.}{Both examples suit great to describe the situation.}
\vocentry{367}{Бросить}{'brosit’}{To throw; to give up}{Как ты мог бросить камень прямо в окно?}{How could you throw a stone right at the window?}
\vocentry{368}{Школа}{sh'kola}{A school}{Эта элитная школа принимает только лучших учеников.}{This elite school accepts only the best students.}
\vocentry{369}{Парень}{'paren’}{A fellow, a guy; a boyfriend}{В общем, он хороший парень, но очень просто выходит из себя.}{He’s a good fellow in general but he loses his temper easily.}
\vocentry{370}{Кровь}{'krov’}{Blood}{Полиция пытается выяснить, чья это кровь.}{The police are trying to find out whose blood it is.}
\vocentry{371}{Двадцать}{'dvatsat’}{Twenty}{Мы приехали на вокзал за двадцать минут до отъезда.}{We arrived at the railway station twenty minutes before the departure.}
\vocentry{372}{Солнце}{'sontse}{The sun}{Солнце всходит на востоке.}{The sun rises in the East.}
\vocentry{373}{Неделя}{n’e'delja}{A week}{Мне нужна неделя, чтобы закончить эксперимент.}{I need a week to finish the experiment.}
\vocentry{374}{Послать}{pas'lat’}{To send}{Мы должны срочно послать за врачом!}{We must send for the doctor urgently!}
\vocentry{375}{Находиться}{naha'ditsa}{To be located}{Я уверен, что этот ресторан просто не может находиться так далеко.}{I’m sure this restaurant just can’t be located so far.}
\vocentry{376}{Ребята}{r’e'bjata}{Guys, young men (often used as vocative)}{Ребята, пошли на каток!}{Guys, let’s go to the skating rink!}
\vocentry{377}{Поставить}{pas'tavit’}{To put, to place (perfective)}{Где я могу поставить свои вещи?}{Where can I put my things?}
\vocentry{378}{Встать}{vstat’}{To stand up}{Он попытался встать, но у него кружилась голова.}{He tried to stand up but he was dizzy.}
\vocentry{379}{Например}{napri'mer}{For example}{Моя дочь не ест многие продукты, например, рыбу.}{My daughter doesn’t eat many foods, fish, for example.}
\vocentry{380}{Шаг}{shag}{Step}{Это решение – очень серьёзный шаг.}{This decision is a very serious step.}
\vocentry{381}{Мужчина}{muzh'china}{A man}{Мой отец – замечательный мужчина.}{My father is a remarkable man.}
\vocentry{382}{Равно}{rav'no}{In combination with ‘всё’ - In any case, It’s all the same}{Зачем мне волноваться? Всё равно это не моя проблема.}{Why should I worry? In any case, it’s not my problem.}
\vocentry{383}{Нос}{nos}{A nose}{Сильный удар сломал ему нос.}{A hard blow broke his nose.}
\vocentry{384}{Мало}{'malə}{Little, few}{Для этой должности у вас слишком мало опыта.}{You have too little experience for this position.}
\vocentry{385}{Внимание}{vni'manie}{Attention}{Иногда внимание лучше любых подарков.}{Sometimes attention is better than any present.}
\vocentry{386}{Капитан}{kapi'tan}{A captain}{Капитан корабля был человеком с отличным чувством юмора.}{The captain of the ship was a man with an excellent sense of humor.}
\vocentry{387}{Ухо}{'uhə}{An ear}{Моя собака глухая на одно ухо.}{My dog is deaf in one ear.}
\vocentry{388}{Туда}{tu'da}{There}{В прошлом году я был в Париже и планирую вернуться туда.}{Last year I was in Paris and I am planning to return there.}
\vocentry{389}{Сюда}{sju'da}{Here, this way}{Как ты попал сюда?}{How did you get here?}
\vocentry{390}{Играть}{ig'rat’}{To play}{Я научился играть на скрипке, когда мне было двенадцать.}{I learned to play the violin when I was twelve.}
\vocentry{391}{Следовать}{'sledavat’}{To follow, to go after}{Я не хочу следовать этим правилам.}{I don’t want to follow these rules.}
\vocentry{392}{Рассказать}{raska'zat’}{To tell (perfective)}{Дети всегда просят меня рассказать им сказку перед сном.}{Children always ask me to tell them a fairy tale before bedtime.}
\vocentry{393}{Великий}{ve'likij}{Great, outstanding}{Пушкин – великий русский поэт.}{Pushkin is a great Russian poet.}
\vocentry{394}{Действительно}{dejst'vitelnə}{Really, actually}{Я действительно ничего не знаю об аварии.}{I really don’t know anything about the accident.}
\vocentry{395}{Слишком}{'slishkəm}{Too}{Современные дети слишком много времени проводят за компьютером.}{Today’s kids spend too much time on the computer.}
\vocentry{396}{Тяжелый}{tja'zhjelyj}{Heavy; difficult, hard}{Я не могу сама нести этот чемодан – он очень тяжёлый.}{I can’t carry this suitcase myself – it’s very heavy.}
\vocentry{397}{Спать}{'spat’}{To sleep}{Она не может спать с включённым светом.}{She can’t sleep with the lights on.}
\vocentry{398}{Оставить}{as'tavit’}{To leave; to abandon}{Как ты мог оставить паспорт дома?}{How could you leave your passport at home?}
\vocentry{399}{Войти}{vaj'ti}{To enter, to come in}{Я не смог войти в комнату – дверь была закрыта.}{I couldn’t enter the room – the door was closed.}
\vocentry{400}{Длинный}{'dlinnyj}{Long}{Я люблю этот старый длинный шарф.}{I love this old long scarf.}
\vocentry{401}{Чувство}{'tchustvə}{A feeling}{Зависть – это плохое чувство.}{Envy is a bad feeling.}
\vocentry{402}{Молчать}{mal'chat’}{To keep silent}{Он попросил меня молчать о нашей ссоре.}{He asked me to keep silent about our argument.}
\vocentry{403}{Рассказывать}{ras'kazyvat’}{To tell}{Мой дедушка умел рассказывать интересные истории.}{My grandfather could tell interesting stories.}
\vocentry{404}{Отвечать}{atve'tchat’}{To answer, to reply}{Могу я не отвечать на этот вопрос?}{May I not answer this question?}
\vocentry{405}{Становиться}{stana'vitsa}{To become; to stand}{Вставать так поздно не должно становиться твоей привычкой.}{Getting up so late shouldn’t become your habit.}
\vocentry{406}{Остановиться}{astana'vitsa}{To stop}{Мальчик понимал, что нельзя так сильно смеяться, но он просто не мог остановиться.}{The boy realized he shouldn’t laugh so much, but he just couldn’t stop.}
\vocentry{407}{Берег}{'bereg}{A bank, a shore}{Берег реки был покрыт густым кустарником.}{The bank of the river was covered with thick bushes.}
\vocentry{408}{Семья}{sem'ja}{A family}{Моя семья всегда меня поддерживает.}{My family always supports me.}
\vocentry{409}{Искать}{is'kat’}{To look for}{Он начал искать работу два месяца назад.}{He started to look for work two months ago.}
\vocentry{410}{Генерал}{gene'ral}{A general}{Мой дядя военный, он уже генерал.}{My uncle is a military man, he’s a general already.}
\vocentry{411}{Момент}{ma'ment}{A moment}{Это уникальный момент в истории государства.}{It’s a unique moment in the history of the state.}
\vocentry{412}{Десять}{'desjat}{Ten}{Мне нужно десять метров шёлка.}{I need ten meters of silk.}
\vocentry{413}{Начать}{na'chat’}{To start, to begin (perfective)}{Я даже не знаю, с чего начать.}{I don’t even know what to start with.}
\vocentry{414}{Следующий}{'sledujuschij}{Next, following}{Каков наш следующий шаг?}{What’s our next step?}
\vocentry{415}{Личный}{'lichnyj}{Personal}{У этого бизнесмена есть личный самолёт.}{This businessman has a personal plane.}
\vocentry{416}{Труд}{trud}{Labor, work}{Тяжёлый труд испортил его здоровье.}{Hard labor spoiled his health.}
\vocentry{417}{Верить}{'verit’}{To believe}{Не следует верить всему, что говорят по телевизору.}{One shouldn’t believe everything they say on TV.}
\vocentry{418}{Группа}{'grupa}{A group}{Первая группа участников оказалась более талантливой.}{The first group of the participants turned out to be more talented.}
\vocentry{419}{Немного}{nem'nogə}{A little, some}{Я могу немного говорить по-русски.}{I can speak Russian a little.}
\vocentry{420}{Впрочем}{'vprotchem}{However}{Впрочем, я не против этого предложения.}{However, I’m not against this proposal.}
\vocentry{421}{Видно}{'vidnə}{Looks like, evidently (colloquial)}{Видно, ты забыл о своём обещании.}{Looks like you forgot about your promise.}
\vocentry{422}{Являться}{jav'ljatsa}{To be (only imperfective); to turn up}{Такое поведение не может являться нормой.}{Such behavior can’t be the norm.}
\vocentry{423}{Муж}{muzh}{A husband}{Её муж – моряк.}{Her husband is a sailor.}
\vocentry{424}{Разве}{'razve}{Really (mostly to express surprise)}{Разве он уже потратил все деньги?}{Has he really spent all the money already?}
\vocentry{425}{Движение}{dvi'zhenije}{Movement, motion}{В балете каждое движение имеет значение.}{Every movement matters in ballet.}
\vocentry{426}{Порядок}{po'rjadək}{Order; sequence}{Пора привести дом в порядок.}{It’s time to put the house in order.}
\vocentry{427}{Ответ}{at'vet}{Answer, reply}{Её ответ удивил всех.}{Her answer surprised everyone.}
\vocentry{428}{Тихо}{ti'ho}{Quietly}{У меня болит горло, поэтому я говорю так тихо.}{I have a sore throat, that’s why I talk so quietly.}
\vocentry{429}{Знакомый}{zna'komyj}{An acquaintance, a friend; familiar}{Мне помог мой знакомый с работы.}{My acquaintance from work helped me.}
\vocentry{430}{Газета}{ga'zeta}{A newspaper}{Каждая газета в стране писала об этом.}{Every newspaper in the country wrote about it.}
\vocentry{431}{Помощь}{'poməstch}{Help}{Моя семья всегда рядом, когда мне нужна помощь.}{My family is always there for me when I need help.}
\vocentry{432}{Сильный}{'sil’nyj}{Strong}{Сильный ветер сдул его шляпу.}{A strong wind blew off his hat.}
\vocentry{433}{Скорый}{'skoryj}{Quick, rapid}{Его скорый отъезд удивил нас всех.}{His quick departure surprised us all.}
\vocentry{434}{Собака}{sa'baka}{A dog}{Эта собака – победитель многих конкурсов.}{This dog is the winner of many competitions.}
\vocentry{435}{Дерево}{'derevə}{A tree}{Мой дедушка посадил это дерево, когда был молодым.}{My grandfather planted this tree when he was young.}
\vocentry{436}{Снег}{sneg}{Snow}{Белый снег искрился на солнце.}{White snow was sparkling in the sun.}
\vocentry{437}{Сон}{son}{Sleep; dream}{Здоровый сон – это главное условие его выздоровления.}{Sound sleep is the main condition for his recovery.}
\vocentry{438}{Смысл}{smysl}{Sense, meaning}{Я не совсем понимаю смысл этой пословицы.}{I don’t quite understand the sense of this proverb.}
\vocentry{439}{Смочь}{smoch’}{To manage, to be able to (perfective, isn’t used in the initial form, in the example below past 3d person singular masculine form is used)}{Наконец, спортсмен смог побить свой собственный рекорд.}{At last, the athlete managed to beat his own record.}
\vocentry{440}{Против}{'protiv}{Against}{Родители были против их отношений.}{The parents were against their relationship.}
\vocentry{441}{Бежать}{be'zhat’}{To run}{Ей пришлось бежать, чтобы успеть на автобус.}{She had to run to catch the bus.}
\vocentry{442}{Двор}{dvor}{A yard}{Мы украсили двор к празднику.}{We decorated the yard for the holiday.}
\vocentry{443}{Форма}{'forma}{A shape, a form}{У этого предмета была странная форма.}{That object had a strange shape.}
\vocentry{444}{Простой}{pras'toj}{Simple}{Это очень простой вопрос, задай мне другой.}{It’s a very simple question, ask me another one.}
\vocentry{445}{Приехать}{pri'jehat’}{To arrive to, to come (perfective)}{Не злись на него – он не мог приехать раньше.}{Don’t be angry with him – he couldn’t arrive earlier.}
\vocentry{446}{Иной}{i'noj}{Different, other, another (more literary style)}{Он так не похож на брата, у него совсем иной характер.}{He’s so much unlike his brother, he has a completely different character.}
\vocentry{447}{Кричать}{kri'tchat’}{To cry, to shout, to scream}{Не нужно кричать – я прекрасно тебя слышу.}{There’s no need to cry – I can hear you perfectly well.}
\vocentry{448}{Возможность}{vaz'mozhnəst’}{An opportunity; a possibility}{Она упустила возможность получить новую работу.}{She missed the opportunity to get a new job.}
\vocentry{449}{Общество}{'obstchestvə}{A society; a community}{Общество не всегда продвигает правильные ценности.}{The society doesn't always promote the right values.}
\vocentry{450}{Зелёный}{ze'l’onyj}{Green}{Считается, что зелёный цвет помогает успокоиться.}{It’s considered that green colors help to calm down.}
\vocentry{451}{Грудь}{grud’}{Chest; breast}{Неожиданный выстрел в грудь убил полицейского.}{A sudden shot in the chest killed the policeman.}
\vocentry{452}{Угол}{'ugəl}{A corner; an angle}{Угол здания был неровным.}{The corner of the building was uneven.}
\vocentry{453}{Открыть}{ot'kryt’}{To open; to discover}{Преступник не смог открыть сейф.}{The criminal didn’t manage to open the safe.}
\vocentry{454}{Происходить}{praisha'dit’}{To take place; to happen}{Подобное не должно происходить в нашей стране.}{Such things shouldn’t happen in our country.}
\vocentry{455}{Ладно}{'ladnə}{All right, okay}{Ладно, я помогу тебе с докладом, но это в последний раз.}{All right, I’ll help you with the report, but it’s the last time.}
\vocentry{456}{Оранжевый}{a'ranzhevyj}{Orange (color)}{Сними этот оранжевый костюм. Ты в нём выглядишь как клоун.}{Take this orange suit off. You look like a clown in it.}
\vocentry{457}{Век}{vek}{Century, age}{Двадцать первый век – это эпоха информационных технологий.}{The twenty first century is the age of information technology.}
\vocentry{458}{Карман}{kar'man}{A pocket}{Карман джинсов не самое надёжное место для ключей.}{A jeans pocket is not the most reliable place for keys.}
\vocentry{459}{Ехать}{'jehat’}{To go (by horse or vehicle), to ride, to drive}{Вы сейчас же должны ехать в больницу.}{You must go to hospital at once.}
\vocentry{460}{Немец}{'nemets}{German, a German man}{Никто не знает, что он немец.}{Nobody knows he’s German.}
\vocentry{461}{Наверное}{na'vernəje}{Probably, likely}{Он, наверное, самый известный русский актёр.}{He’s probably the most famous Russian actor.}
\vocentry{462}{Губа}{gu'ba}{A lip}{Верхняя губа ребенка распухла от укуса пчелы.}{The child’s upper lip is swollen from a bee sting.}
\vocentry{463}{Дядя}{'djadja}{An uncle; a man (colloquial)}{Мой дядя профессиональный фотограф.}{My uncle is a professional photographer.}
\vocentry{464}{Приходить}{priha'dit’}{To come; to arrive at}{Я обещаю приходить каждый день и навещать тебя в больнице.}{I promise to come every day and visit you in the hospital.}
\vocentry{465}{Часто}{'tchastə}{Often}{Мы с женой часто вспоминаем день нашего знакомства.}{My wife and I often recollect the day we met.}
\vocentry{466}{Домой}{da'moj}{Home (direction, where to)}{Темнеет, мне пора домой.}{It’s getting dark, it’s time for me to go home.}
\vocentry{467}{Огонь}{a'gon’}{Fire}{Твои слова только добавляют масла в огонь.}{Your words are only adding fuel to the fire.}
\vocentry{468}{Писатель}{pi'satel’}{A writer}{Кто твой любимый американский писатель?}{Who’s your favorite American writer?}
\vocentry{469}{Армия}{'armija}{The army}{Армия должна защищать население страны.}{The army must protect the population of the country.}
\vocentry{470}{Состояние}{sasta'janije}{State, condition; fortune}{Состояние этой машины оставляет желать лучшего.}{The state of this car leaves much to be desired.}
\vocentry{471}{Зуб}{zub}{A tooth}{Стоматолог аккуратно удалил гнилой зуб.}{The dentist carefully removed the decayed tooth.}
\vocentry{472}{Очередь}{'otchered’}{A queue}{Перед Эрмитажем всегда большая очередь.}{There’s always a long queue in front of the Hermitage.}
\vocentry{473}{Кой}{koj}{In the expression ‘на кой чёрт’ - ‘Why the hell?’}{На кой чёрт ты купил все эти вещи?}{Why the hell have you bought all these things?}
\vocentry{474}{Подняться}{pad'n’atsatsa}{To rise, to go up, to get up}{Как цены могли подняться так быстро?}{How could the prices rise so quickly?}
\vocentry{475}{Камень}{'kamen’}{A stone; a rock}{Вчера из музея был украден драгоценный камень.}{Yesterday a precious stone was stolen from the museum.}
\vocentry{476}{Гость}{gost’}{A guest}{Расслабься и получай удовольствие – ты наш гость.}{Relax and enjoy – you’re our guest.}
\vocentry{477}{Показаться}{paka'zatsa}{To seem, to appear}{Её поведение может показаться странным, но ты привыкнешь к этому.}{Her behavior may seem strange, but you’ll get used to it.}
\vocentry{478}{Ветер}{'veter}{Wind}{Вчерашний сильный ветер нанес ущерб многим домам.}{Yesterday’s strong wind did harm to many houses.}
\vocentry{479}{Собираться}{sabi'ratsa}{To gather; to be about to}{Мы любим собираться у родителей по праздникам.}{We like to gather at our parents’ on holidays.}
\vocentry{480}{Попасть}{pa'past’}{To get to; to hit (a target) (perfective)}{Как ты смог попасть сюда без билета?}{How did you manage to get here without a ticket?}
\vocentry{481}{Принять}{pri'njat’}{To take; to admit (perfective)}{Мы должны принять срочные меры, пока ещё не поздно.}{We must take urgent measures before it’s too late.}
\vocentry{482}{Сначала}{sna'tchala}{At first}{Сначала я не был уверен в своём решении.}{At first I wasn’t sure about my decision.}
\vocentry{483}{Либо}{'libə}{Either … or, or}{Он либо действительно ничего не знает, либо притворяется.}{He either really doesn’t know anything or is pretending.}
\vocentry{484}{Поехать}{pa'jehat’}{To go (by horse or vehicle), to ride, to drive (perfective)}{Извини, я не смогу поехать с тобой за город.}{Sorry, I won’t be able to go to the country with you.}
\vocentry{485}{Услышать}{us'lyshat’}{To hear (perfective)}{Я бы так хотела услышать её пение!}{I would so much like to hear her singing!}
\vocentry{486}{Уметь}{u'met’}{To be able to, to know how to}{Ты должен уметь прощать, если хочешь сохранить свою семью.}{You should be able to forgive if you want to keep your family.}
\vocentry{487}{Случиться}{slu'tchitsa}{To happen, to occur}{Не волнуйся так сильно! Что могло случиться?}{Don’t worry so much! What could have happened?}
\vocentry{488}{Странный}{'strannyj}{Strange}{Что это за странный запах?}{What kind of strange smell is it?}
\vocentry{489}{Единственный}{je'dinstvennyj}{Only, sole}{Он единственный ребёнок в семье.}{He’s an only child in the family.}
\vocentry{490}{Рота}{'rota}{Company (military)}{Одна рота не может сражаться против целой армии.}{One company can’t fight against the whole army.}
\vocentry{491}{Закон}{za'kon}{Law}{Закон один для всех.}{The law is the same for everyone.}
\vocentry{492}{Короткий}{ka'rotkij}{Short}{Этот курс короткий, но очень полезный.}{This course is short but very useful.}
\vocentry{493}{Море}{'more}{Sea}{Мёртвое море – самое солёное в мире.}{The Dead Sea is the saltiest in the world.}
\vocentry{494}{Добрый}{'dobryj}{Kind}{Мой друг очень добрый и всегда готов помочь.}{My friend is very kind and is always ready to help.}
\vocentry{495}{Тёмный}{'tjemnyj}{Dark}{Большинство детей не любят тёмный шоколад.}{Most children don’t like dark chocolate.}
\vocentry{496}{Гора}{ga'ra}{A mountain}{Эта гора выглядит очень живописно.}{This mountain looks very picturesque.}
\vocentry{497}{Врач}{vratch}{A doctor, a physician}{Наш семейный врач – очень ответственный человек.}{Our family doctor is a very responsible man.}
\vocentry{498}{Край}{kraj}{An edge; a land}{Край ножа был очень острым.}{The edge of the knife was very sharp.}
\vocentry{499}{Стараться}{sta'ratsa}{To try, to attempt}{Я буду стараться изо всех сил.}{I’ll try my best.}
\vocentry{500}{Лучший}{'lutchij}{The best}{Он лучший игрок в команде.}{He’s the best player in the team.}
\end{multicols}
